The deterministic neural model is the non-autonomous Neural ODE
\citep{chen2018neuralode}
\begin{equation}
  \frac{d z}{d\tau} = f_\theta(\tau,z),
  \qquad
  z(0)=z_0.
\end{equation}

The fitted solution is a single population curve, not one curve per trial. The
three trials are treated as replicate observations of the same underlying
trajectory. On the original response scale the initial condition is
\begin{equation}
    y_0 = y(t_0)
    =
        \frac{1}{J}\sum_{j=1}^J y_0^{[j]},
\end{equation}
%
the mean of the first observed values across trials. The corresponding
normalized initial condition is
$$
    z_0 = \frac{y_0-\bar y}{s_y}.
$$
After solving the normalized ODE, the fitted value on the original scale is
\begin{equation}
    \label{eqn:inverze_normalization}
    \begin{aligned}
        \widehat y_i
            &=
                \widehat y(t_i)
        \\
            &=
            \bar y +
            s_y \cdot \widehat z(\tau_i;\theta)
    \end{aligned}
\end{equation}
The same value $\widehat y_i$ is compared with every trial observation
$y_i^{[j]}$ at time $t_i$.

\paragraph{Network Architecture}
The right-hand side $f_\theta$ is a fully connected neural network with input
$(\tau,z)\in\mathbb R^2$, two hidden layers, 32 hidden units per layer, and
hyperbolic tangent activation:

\begin{equation}
    \label{eqn:activation_functions}
    \begin{aligned}
        h_1 &=
            \tanh(W_1 x+b_1),
        \\
        h_{\ell}
            &=
                \tanh(W_{\ell}h_{\ell-1}+b_{\ell}),
        \\
        f_\theta(x)
            &=
                W_o h_2+b_o
    \end{aligned}
\end{equation}
where $x=(\tau,z)^T$. The network has 1185 trainable weights and biases.
This count comes from the three affine layers:
$$
\begin{aligned}
  \hbox{input to hidden 1:}&\quad 2\cdot 32+32=96,\\
  \hbox{hidden 1 to hidden 2:}&\quad 32\cdot 32+32=1056,\\
  \hbox{hidden 2 to output:}&\quad 32\cdot 1+1=33.
\end{aligned}
$$
Thus $96+1056+33=1185$. The architecture was chosen as a compact smooth
function approximator for a scalar time-dependent vector field. The input has
only two coordinates, time and state, and the data set has only 45 scalar
observations, so a very large network would make the comparison dominated by
overparameterization. Two hidden layers with 32 units give substantially more
flexibility than the two-parameter logistic curve while remaining small enough
to train stably with weight decay. The $\tanh$ nonlinearity gives a smooth
right-hand side, which is useful for ODE/SDE integration, and sigmoidal neural
networks are classical universal approximators \citep{cybenko1989approximation}.
The same architecture is used for the Neural ODE and for each Neural SDE
coefficient so that differences between the deterministic and stochastic neural
models are not caused by different network sizes.

The ODE was solved at the observed normalized times using RK4 with step size
$0.01$. The empirical loss was
$$
  \mathcal L_{\mathrm{ODE}}(\theta)
  =
  \frac{1}{n}
  \sum_{(i,j)\in\Omega}
  \left(z_i^{[j]}-\widehat z_i(\theta)\right)^2.
$$
Equivalently, in original response units this loss fits one curve
$\widehat y_i$ to all three replicate measurements at the same time. The
minimizer is therefore encouraged to pass through the center of the replicate
cloud, while the ODE structure enforces a smooth dynamical trajectory.
Training used Adam for 1000 epochs with learning rate $10^{-2}$ and weight
decay $10^{-4}$ \citep{kingma2015adam}.
